\section{Proposed Method}
\label{sec:methodology}

% Follows G_paper_and_latex/paper_outline.md (Section 3) and
% E_research_design/architecture_candidates.md.
% All formulations, dimensions, and blocks must represent the actual implemented system.

\subsection{Problem Formulation}
\label{sec:formulation}
% TODO: Define input image space X \in R^{H x W x 3}, bounding box parameters (x, y, w, h),
% class label set C, and optimization objective.
\noindent [TODO: Formally define the 2D multi-class object detection task in warehouse scenes: input resolution $H \times W \times 3$, ground truth bounding box format, class set $\mathcal{C}$ derived from the validated dataset taxonomy, and overall optimization objectives. Reference \texttt{E\_research\_design/research\_questions\_and\_hypotheses.md}].

\subsection{Overall Architecture Overview}
\label{sec:arch_overview}
% Follows yolo11s_ss2d_architecture_implementation.md and yolo_mamba_research_cases.md
The proposed hybrid architecture anchors on the verified lightweight \textbf{YOLO11s} detector. To achieve deployable long-range spatial context reasoning without suffering the quadratic memory scaling $\mathcal{O}(N^2)$ of vision transformers or excessive Memory Access Costs (MAC) at large resolutions, the model enforces a strict stage-wise separation:
\begin{enumerate}
    \item \textbf{High-Resolution Stages ($P_1$--$P_3$):} Retained as standard local convolutions to preserve edge gradients and fine details with optimal cache locality.
    \item \textbf{Deep-Stage Context Modeling ($P_5$ / $P_4$):} Selectively replaces the deepest Pointwise Spatial Attention block (\texttt{C2PSA}) with an input/output shape-compatible 2D Selective Scan (\textbf{SS2D}) block accepting $[B, C, H, W]$.
    \item \textbf{Feature Neck \& Head:} Preserves the original unmodified YOLO11s FPN/PAN neck and decoupled anchor-free \texttt{Detect} head to maintain baseline stability and enable clean, unconfounded ablation isolation.
\end{enumerate}

\subsection{CNN Local Feature Extraction}
\label{sec:cnn_features}
% TODO: Specify the exact CNN backbone blocks used (e.g., CSP/C2f/C3k2 blocks),
% layer channels, strides, and receptive field characteristics.
\noindent [TODO: Specify the exact CNN backbone architecture, layer dimensions, feature strides ($P_3, P_4, P_5$), and inductive bias for local feature extraction].

\subsection{Mamba/State-Space Contextual Block}
\label{sec:mamba_integration}
% TODO: 2D-to-1D sequence flattening, scanning mechanisms (e.g., 2-way or 4-way selective scan),
% state dimension, projection matrices, and computational implications.
\noindent [TODO: Define the mathematical formulation of the Mamba state-space block, including sequence conversion from 2D feature maps, scanning directions, discretization parameters, and linear state updates. Detail exact tensor dimensions and channel projections; see architecture candidates documentation].

\subsection{Multi-Scale Feature Fusion}
\label{sec:fusion}
% TODO: Mathematical definition of the feature fusion mechanism (e.g., residual addition,
% concatenation with projection, or learned gating).
\noindent [TODO: Present the formal equations for fusing local CNN representations and global Mamba contextual representations. Specify channel alignment and projection operations; reference architecture candidates documentation].

\subsection{Detection Head and Loss Functions}
\label{sec:head_loss}
% TODO: Anchor-free / anchor-based head formulation, classification loss (BCE/Focal),
% bounding-box regression loss (CIoU/GIoU/DFL).
\noindent [TODO: Detail the detection head architecture and multi-task loss formulation (classification loss, box regression loss, and distribution focal loss). Specify hyperparameters and loss weighting coefficients].

\subsection{Computational Complexity Analysis}
\label{sec:complexity}
% TODO: Analytical parameter count, FLOPs/MACs calculation, and comparison of asymptotic
% complexity: O(N) for SSM vs O(N^2) for Multi-Head Self-Attention.
\noindent [TODO: Provide analytical complexity formulas comparing the hybrid block with standard self-attention mechanisms, detailing parameter footprint, FLOPs/MACs at standard input resolution ($640 \times 640$), and theoretical memory bounds. See \texttt{E\_research\_design/architecture\_candidates.md} (Section 4)].

\subsection{Training and Inference Workflow}
\label{sec:training_workflow}
% TODO: End-to-end training pipeline, optimization schedule, checkpoint selection,
% and deployment conversion steps (e.g., export to ONNX/TensorRT).
\noindent [TODO: Document the step-by-step training pipeline, convergence criteria, and export workflow for edge inference runtimes. Reference \texttt{E\_research\_design/baseline\_plan.md}].
